% ============================================================================
% CAPÍTULO 3 — METODOLOGIA
% ============================================================================
\chapter{Metodologia}
\label{cap:metodologia}

\section{Delineamento geral}

Este trabalho utiliza uma abordagem de simulação: agentes de inteligência
artificial baseados em modelos de linguagem de grande escala (LLM) recebem,
via instrução de sistema (\textit{system prompt}), uma personalidade
induzida segundo o modelo Big Five, e em seguida respondem a dois
instrumentos: um de autorrelato de personalidade (checagem de manipulação)
e o Índice de Perfil do Investidor (IPA, variável dependente).

Trata-se de uma pesquisa \textbf{quantitativa}, pois a relação entre
personalidade induzida e escore do IPA é testada por meios estatísticos
(correlação de Spearman e regressão OLS, seção \ref{sec:analise-estatistica}),
não por análise qualitativa das respostas dos agentes. É também
\textbf{exploratória}, pois não foi encontrado um estudo anterior que
relacione personalidade induzida em agentes de IA a um questionário de
perfil de investidor (\textit{suitability}): os estudos mais próximos com
LLMs \cite{borman2024,hartley2025} medem o risco por tarefas de escolha,
pela avaliação de loterias e por questionários comportamentais, não por um questionário de
\textit{suitability}, e as hipóteses de referência deste trabalho
(Capítulo \ref{cap:referencial}) vêm de estudos com sujeitos humanos,
adaptadas por analogia.

O método é o de \textbf{simulação computacional controlada}:
diferentemente de um levantamento ou experimento com sujeitos humanos, o
nível de personalidade solicitado a cada agente é conhecido por construção
(definido pelo sorteio que gera o \textit{prompt}), o que elimina o erro
de medida da variável independente presente em estudos com autorrelato
humano. Em contrapartida, a validade das conclusões depende de o LLM realmente
incorporar, em sua geração de texto, a personalidade descrita no
\textit{prompt}; daí a necessidade da checagem de manipulação (seção
\ref{sec:bfi44}) e da verificação do desenho de indução do Experimento B (seção
\ref{sec:exp-a-b}).

\section{Instrumentos}

\subsection{Indução de personalidade (Baixo/Médio/Alto)}

Para cada um dos cinco traços do Big Five foram escritas três descrições
curtas (Baixo, Médio, Alto), baseadas no conteúdo dos itens do BFI-10
(seção \ref{sec:bfi10-piloto}). Essas
descrições são inseridas no \textit{system prompt} de cada agente para
induzir a personalidade (seção \ref{sec:prompts}). O Quadro
\ref{quadro:perfis-nivel} reproduz as 15 descrições.

\begin{quadro}[htbp]
    \caption{Descrições usadas para induzir cada nível de cada traço}
    \label{quadro:perfis-nivel}
    \centering
    \small
    \begin{tabular}{lp{12cm}}
        \toprule
        Nível & Descrição inserida no \textit{prompt} \\
        \midrule
        \multicolumn{2}{l}{\textbf{Extroversão}} \\
        Baixo & É uma pessoa reservada, prefere ambientes tranquilos e evita ser o centro das atenções. \\
        Médio & É uma pessoa equilibrada entre momentos sociais e momentos sozinha, nem muito reservada nem muito extrovertida. \\
        Alto & É uma pessoa sociável e extrovertida, gosta de estar com outras pessoas e de ambientes movimentados. \\
        \addlinespace
        \multicolumn{2}{l}{\textbf{Amabilidade}} \\
        Baixo & Tende a ser crítica com os outros, aponta defeitos com facilidade e confia pouco nas pessoas. \\
        Médio & Confia nas pessoas moderadamente: às vezes é crítica, às vezes é tolerante. \\
        Alto & Geralmente confia nas pessoas, é gentil e raramente é crítica ou aponta defeitos nos outros. \\
        \addlinespace
        \multicolumn{2}{l}{\textbf{Consciência}} \\
        Baixo & Tende a ser preguiçosa, desorganizada e costuma deixar tarefas inacabadas. \\
        Médio & É razoavelmente organizada, mas às vezes deixa tarefas para depois. \\
        Alto & É organizada e disciplinada, insiste até concluir qualquer tarefa ou trabalho. \\
        \addlinespace
        \multicolumn{2}{l}{\textbf{Neuroticismo}} \\
        Baixo & É uma pessoa relaxada, que controla bem o estresse e raramente fica nervosa. \\
        Médio & Mantém o equilíbrio emocional na maior parte do tempo, mas ocasionalmente fica ansiosa diante de problemas. \\
        Alto & Fica nervosa com facilidade e tem dificuldade para controlar o estresse e a ansiedade. \\
        \addlinespace
        \multicolumn{2}{l}{\textbf{Abertura à Experiência}} \\
        Baixo & Tem poucos interesses artísticos e pouca imaginação, prefere o familiar e o convencional. \\
        Médio & Tem alguns interesses artísticos e está moderadamente aberta a novas ideias e experiências. \\
        Alto & Tem uma imaginação fértil, muitos interesses artísticos e intelectuais, e gosta de experimentar coisas novas. \\
        \bottomrule
    \end{tabular}
    \legend{Fonte: elaborado pela autora.}
\end{quadro}

\subsection{Questionário de autorrelato: checagem de manipulação}
\label{sec:bfi44}

O instrumento utilizado como checagem de manipulação na análise principal
deste trabalho é o \textbf{BFI-44} \cite{johndonahuekentle1991}, versão
completa de 44 itens do \textit{Big Five Inventory}, em uma versão em
português que segue a ordem e o conteúdo dos 44 itens do instrumento
original, com a mesma escala de concordância de 1
(discordo totalmente) a 5 (concordo totalmente) usada no BFI-10 (seção
\ref{sec:bfi10-piloto}). A validação psicométrica do BFI-44 para o
Brasil foi conduzida por \citeonline{andrade2008}, com uma adaptação cuja
ordem e redação dos itens diferem em parte das usadas aqui; por isso, as
propriedades psicométricas relatadas por esse autor servem como
referência geral para o instrumento, e não como validação da versão
exata aplicada aos agentes. Cada um dos cinco traços é medido por um número
diferente de itens: Extroversão e Neuroticismo com 8 cada, Abertura à
Experiência com 10, e Amabilidade e Consciência com 9 cada, com itens
diretos e reversos em proporção variável por traço; a pontuação de cada traço é a média
dos seus itens, já com a reversão aplicada quando necessário (itens
reversos: escore final = 6 menos a resposta, mesma regra do BFI-10).

A Tabela \ref{tab:bfi44-resumo} resume o número de itens por traço, e o
Quadro \ref{quadro:bfi44-itens} reproduz as 44 afirmações, na forma em
que foram apresentadas aos agentes. As instruções e a escala de resposta
foram:

\begin{itemize}
    \item \textbf{Instruções}: ``As declarações a seguir dizem respeito à
    sua percepção sobre si mesmo em uma variedade de situações. Indique o
    número que melhor expressa sua opinião sobre você mesmo em cada
    afirmação.'' Frase-base: ``Eu me vejo como alguém que...'';
    \item \textbf{Escala de resposta}: 1 = Discordo totalmente; 2 =
    Discordo em parte; 3 = Nem concordo nem discordo; 4 = Concordo em
    parte; 5 = Concordo totalmente.
\end{itemize}

\begin{quadro}[htbp]
    \caption{Itens do BFI-44}
    \label{quadro:bfi44-itens}
    \centering
    \footnotesize
    \begin{tabular}{rp{7.6cm}lc}
        \toprule
        Item & Afirmação (``Eu me vejo como alguém que...'') & Traço & Escore \\
        \midrule
        1 & Gosta de conversar, é comunicativa. & Extroversão & D \\
        2 & Tende a criticar os outros. & Amabilidade & R \\
        3 & É minuciosa e detalhista no trabalho. & Consciência & D \\
        4 & Depressiva, triste. & Neuroticismo & D \\
        5 & É original, tem ideias novas. & Abertura à Experiência & D \\
        6 & É reservada. & Extroversão & R \\
        7 & É generosa e não é egoísta com outras pessoas. & Amabilidade & D \\
        8 & Pode ser desleixada para fazer as coisas. & Consciência & R \\
        9 & É tranquila, lida bem com estresse. & Neuroticismo & R \\
        10 & Se interessa por áreas diferentes do conhecimento. & Abertura à Experiência & D \\
        11 & É cheia de energia. & Extroversão & D \\
        12 & Inicia bate-boca com outros. & Amabilidade & R \\
        13 & É confiável no trabalho. & Consciência & D \\
        14 & Pode ser tensa. & Neuroticismo & D \\
        15 & É inovadora, pensa profundamente nas coisas. & Abertura à Experiência & D \\
        16 & Gera muito entusiasmo. & Extroversão & D \\
        17 & Desculpa, perdoa os outros. & Amabilidade & D \\
        18 & Tende a ser desorganizada. & Consciência & R \\
        19 & Se preocupa muito, em excesso. & Neuroticismo & D \\
        20 & Tem uma imaginação fértil. & Abertura à Experiência & D \\
        21 & Tende a ser quieta. & Extroversão & R \\
        22 & Geralmente confia, acredita nos outros. & Amabilidade & D \\
        23 & Tende a ser preguiçosa. & Consciência & R \\
        24 & É emocionalmente estável, não se perturba facilmente. & Neuroticismo & R \\
        25 & É inventiva. & Abertura à Experiência & D \\
        26 & É assertiva, não tem medo de expressar o que sente. & Extroversão & D \\
        27 & Às vezes é indiferente com os outros. & Amabilidade & R \\
        28 & Persevera até concluir as tarefas. & Consciência & D \\
        29 & É temperamental e instável emocionalmente. & Neuroticismo & D \\
        30 & Valoriza experiências artísticas e estéticas. & Abertura à Experiência & D \\
        31 & Às vezes é tímida, inibida. & Extroversão & R \\
        32 & É boa e atenciosa com quase todo mundo. & Amabilidade & D \\
        33 & Faz as coisas com eficiência. & Consciência & D \\
        34 & Se mantém calma em situações tensas. & Neuroticismo & R \\
        35 & Gosta de rotina. & Abertura à Experiência & R \\
        36 & É extrovertida e sociável. & Extroversão & D \\
        37 & Às vezes é grosseira com outras pessoas. & Amabilidade & R \\
        38 & Cumpre, finaliza os planos que faz. & Consciência & D \\
        39 & Fica nervosa facilmente. & Neuroticismo & D \\
        40 & Gosta de refletir, jogar com as ideias. & Abertura à Experiência & D \\
        41 & Tem poucos interesses artísticos. & Abertura à Experiência & R \\
        42 & Gosta de cooperar com outros. & Amabilidade & D \\
        43 & Se distrai facilmente. & Consciência & R \\
        44 & É sofisticada em arte, música ou literatura. & Abertura à Experiência & D \\
        \bottomrule
    \end{tabular}
    \legend{Fonte: itens conforme aplicados neste trabalho, versão em português do instrumento de \citeonline{johndonahuekentle1991}. D = escore direto; R = escore reverso.}
\end{quadro}

\begin{table}[htb]
    \caption{Número de itens do BFI-44, por traço}
    \label{tab:bfi44-resumo}
    \centering
    \begin{tabular}{lccc}
        \toprule
        Traço & Itens (total) & Diretos & Reversos \\
        \midrule
        Abertura à Experiência & 10 & 8 & 2 \\
        Amabilidade             & 9  & 5 & 4 \\
        Consciência             & 9  & 5 & 4 \\
        Extroversão             & 8  & 5 & 3 \\
        Neuroticismo            & 8  & 5 & 3 \\
        \bottomrule
    \end{tabular}
    \legend{Fonte: elaborado pela autora.}
\end{table}

Por ter mais itens por traço que o BFI-10 (2 por traço), o BFI-44 produz
uma média menos sensível ao ruído de um único item. Foi esse o principal
motivo para substituir o BFI-10, usado no estudo piloto, pelo BFI-44 na
análise principal: uma checagem de manipulação mais granular permite
distinguir com mais confiança entre ``o traço não foi incorporado pelo
agente'' e ``o traço foi incorporado, mas um ou dois itens específicos tiveram
resposta ruidosa'' (ver seção \ref{sec:limitacoes-pipeline} sobre o viés de
itens de pontuação reversa, que afeta mais um instrumento curto do que
um longo).

\subsection{Estudo piloto: BFI-10}
\label{sec:bfi10-piloto}

Antes da análise principal com o BFI-44 (seção \ref{sec:bfi44}), foi
conduzido um estudo piloto com o \textbf{BFI-10} \cite{rammstedtjohn2007},
versão reduzida de 10 itens do \textit{Big Five Inventory}, na tradução
em português reproduzida no Anexo 3 (material suplementar S2
Questionnaire) de \citeonline{debortoli2019}. Cada um dos cinco traços do Big Five
é medido por 2 itens, um com escore direto e outro com escore reverso, e
a pontuação da dimensão é a média dos seus 2 itens (mesma regra de
reversão do BFI-44).

O Quadro \ref{quadro:bfi10-itens} reproduz os 10 itens, com o traço
medido por cada um. As instruções e a escala de resposta são as mesmas
do BFI-44 (seção \ref{sec:bfi44}).

\begin{quadro}[htbp]
    \caption{Itens do BFI-10}
    \label{quadro:bfi10-itens}
    \centering
    \small
    \begin{tabular}{rp{7.2cm}lc}
        \toprule
        Item & Afirmação (``Eu me vejo como alguém que...'') & Traço & Escore \\
        \midrule
        1  & É reservado. & Extroversão & R \\
        2  & Geralmente confia nas pessoas. & Amabilidade & D \\
        3  & Tende a ser preguiçoso. & Consciência & R \\
        4  & É relaxado, controla bem o estresse. & Neuroticismo & R \\
        5  & Tem poucos interesses artísticos. & Abertura à Experiência & R \\
        6  & É sociável, extrovertido. & Extroversão & D \\
        7  & Tende a ser crítico com os outros (encontrar defeitos). & Amabilidade & R \\
        8  & Insiste até concluir a tarefa ou o trabalho. & Consciência & D \\
        9  & Fica nervoso facilmente. & Neuroticismo & D \\
        10 & Tem uma imaginação fértil. & Abertura à Experiência & D \\
        \bottomrule
    \end{tabular}
    \legend{Fonte: adaptado de \citeonline{debortoli2019}, Anexo 3, que traduz \citeonline{rammstedtjohn2007}. D = escore direto; R = escore reverso.}
\end{quadro}

O BFI-10 foi mantido deliberadamente como instrumento fixo em todos os
experimentos de comparação entre modelos de linguagem (seção
\ref{sec:selecao-modelo}): \texttt{llama3.2:1b}, \texttt{llama3.2:3b},
\texttt{gemma3:4b} e \texttt{llama3.1:8b}. Manter o mesmo instrumento de checagem
de manipulação em todos os modelos testados é o que torna válida a
comparação entre eles; usar instrumentos diferentes para modelos diferentes
confundiria efeito do modelo com efeito do instrumento. Os resultados do
estudo piloto (N=300) são
reportados na seção \ref{sec:resultados-piloto}, e permanecem relevantes
mesmo após a migração para o BFI-44: foram eles que motivaram a escolha do
\texttt{llama3.1:8b} como modelo principal (seção
\ref{sec:selecao-modelo}) e produziram o primeiro indício do achado
central deste trabalho (Abertura à Experiência), posteriormente
confirmado com o instrumento mais granular.

\subsection{Experimentos A e B: indução simultânea vs. isolada de traços}
\label{sec:exp-a-b}

Além da troca de instrumento (do BFI-10 para o BFI-44), dois desenhos de
indução foram testados com o BFI-44, mantendo o mesmo modelo
(\texttt{llama3.1:8b}) e N=300 em cada um:

\begin{itemize}
    \item \textbf{Experimento A (multi-traço)}: mesmo desenho do estudo
    piloto, em que os cinco traços variam Baixo/Médio/Alto de forma
    independente e simultânea em cada agente (seção
    \ref{sec:desenho-fatorial}). É o conjunto de dados da \textbf{análise
    principal} deste trabalho, reportado no Capítulo \ref{cap:resultados};
    \item \textbf{Experimento B (traço dominante)}: cada agente tem
    apenas um dos cinco traços sorteado para variar Baixo/Médio/Alto; os
    outros quatro permanecem fixos em ``Médio'' no \textit{system prompt}.
\end{itemize}

O Experimento B não substitui o A: é uma verificação do desenho de indução,
motivada pela possibilidade de que induzir os cinco traços
simultaneamente no mesmo \textit{prompt} dilua ou confunda o sinal de
cada traço individual, em comparação com a manipulação de apenas um traço
por vez. Se as correlações entre nível induzido e autorrelato (checagem
de manipulação) forem comparáveis nos dois desenhos, haverá evidência de
que a indução simultânea não compromete a qualidade da manipulação. A
indução simultânea é também o desenho mais próximo da personalidade
humana, na qual os cinco traços coexistem sempre em algum nível, e
corresponde à forma como os traços são medidos em
\citeonline{debortoli2019} e \citeonline{jiangpengyan2024}. Os resultados
dessa comparação estão na seção \ref{sec:resultados-exp-b} e são
interpretados no Capítulo \ref{cap:discussao}.

\subsection{IPA: perfil de investidor (variável dependente)}

O IPA é um instrumento de 8 perguntas objetivas (alternativas a--d),
adaptado do questionário de Análise de Perfil do Investidor do Banco do
Brasil, alinhado às diretrizes da ANBIMA e da IOSCO
(\textit{suitability}), conforme reproduzido no Anexo 5 de
\citeonline{debortoli2019}; o nome IPA (Índice de Perfil do Investidor)
é o adotado neste trabalho para esse instrumento. O Quadro
\ref{quadro:ipa-perguntas} reproduz as perguntas e as alternativas.

\begin{quadro}[htbp]
    \caption{Perguntas do IPA e alternativas}
    \label{quadro:ipa-perguntas}
    \centering
    \small
    \begin{tabular}{cp{13.4cm}}
        \toprule
        \multicolumn{2}{p{14.4cm}}{\textbf{1. Indique sua faixa etária.}} \\
        & a) Abaixo de 25 anos \quad b) De 25 a 40 anos \quad c) De 41 a 55 anos \quad d) Acima de 56 anos \\
        \addlinespace
        \multicolumn{2}{p{14.4cm}}{\textbf{2. Se você realizasse um investimento financeiro hoje (ex: fundos de investimento, poupança, CDI), qual o percentual do total de seu patrimônio que você aplicaria em investimentos?}} \\
        & a) Até 25\% \quad b) Entre 26\% e 50\% \quad c) Entre 51\% e 75\% \quad d) Mais de 75\% \\
        \addlinespace
        \multicolumn{2}{p{14.4cm}}{\textbf{3. Quando você realiza alguma aplicação (investimento financeiro), qual o prazo médio dos investimentos?}} \\
        & a) Indefinidamente, pois não tenho planos para utilização \quad b) No máximo 5 anos \quad c) No máximo 2 anos \quad d) No máximo 1 ano \\
        \addlinespace
        \multicolumn{2}{p{14.4cm}}{\textbf{4. Qual o principal objetivo de seus investimentos?}} \\
        & a) Obter rentabilidade superior à proporcionada por investimentos tradicionais de renda fixa, mesmo assumindo risco de possíveis perdas \quad b) Diversificar os investimentos \quad c) Formar uma reserva financeira para utilização futura \quad d) Preservar o patrimônio \\
        \addlinespace
        \multicolumn{2}{p{14.4cm}}{\textbf{5. Qual a sua melhor referência de rentabilidade?}} \\
        & a) Índices das Bolsas de Valores \quad b) Dólar \quad c) CDI \quad d) Poupança \\
        \addlinespace
        \multicolumn{2}{p{14.4cm}}{\textbf{6. Quais são as aplicações financeiras em que você tem maior conhecimento?}} \\
        & a) Fundo de Ações, Ações ou Derivativos \quad b) Fundos Multimercados, Cambiais, Dívida Externa ou Renda Fixa \quad c) CDB, Poupança ou Fundos Referenciados DI e Curto Prazo \quad d) Não tenho conhecimento \\
        \addlinespace
        \multicolumn{2}{p{14.4cm}}{\textbf{7. Você já investiu em ações ou Fundo de Ações?}} \\
        & a) Sim, pois a possibilidade de maiores rentabilidades me atraem muito \quad b) Sim, mas com receio \quad c) Não, mas poderia investir num momento oportuno \quad d) Não e não pretendo investir, pois não me agrada a ideia de estar sujeito a rentabilidade negativa \\
        \addlinespace
        \multicolumn{2}{p{14.4cm}}{\textbf{8. Caso a sua aplicação tivesse uma perda no curto prazo, qual seria o percentual aceitável?}} \\
        & a) Acima de 15\% \quad b) Até 15\% \quad c) Até 5\% \quad d) Não aceitaria perda \\
        \bottomrule
    \end{tabular}
    \legend{Fonte: adaptado de \citeonline{debortoli2019}, Anexo 5 (questionário de Análise de Perfil do Investidor do Banco do Brasil, alinhado à ANBIMA e à IOSCO).}
\end{quadro}

A pergunta 1 (faixa etária) foi atribuída aleatoriamente a cada agente, e
não respondida pelo modelo, por ser um dado demográfico inaplicável a
agentes de IA (ver seção \ref{sec:limitacoes-pipeline}).

\subsubsection{Regra de pontuação do IPA}
\label{sec:ipa-pontuacao}

O artigo de \citeonline{debortoli2019} não publica os pesos atribuídos a
cada alternativa nem os cortes de classificação de perfil
(conservador/moderado/arrojado) usados pelos autores. A regra de
pontuação descrita nesta seção (Tabelas \ref{tab:ipa-pontuacao-alternativas}
e \ref{tab:ipa-classificacao}) é, portanto, uma \textbf{definição própria
deste trabalho}, definida pelo critério mais simples possível para viabilizar o cálculo
automático do perfil de cada agente.

Cada alternativa recebe de 1 a 4 pontos, em que 1 corresponde à
alternativa mais associada a um perfil conservador e 4 à mais associada a
um perfil arrojado. Em 7 das 8 perguntas, a pontuação segue a ordem das
alternativas (a = 4, b = 3, c = 2, d = 1). A exceção é a pergunta 2, em
que a ordem é invertida, pois aplicar uma fração maior do patrimônio em
investimentos indica maior tolerância ao risco, não menor (Tabela
\ref{tab:ipa-pontuacao-alternativas}).

\begin{table}[htb]
    \caption{Pontuação de cada alternativa do IPA, por pergunta}
    \label{tab:ipa-pontuacao-alternativas}
    \centering
    \begin{tabular}{cccccc}
        \toprule
        Pergunta & a & b & c & d & Regra \\
        \midrule
        1, 3, 4, 5, 6, 7, 8 & 4 & 3 & 2 & 1 & Geral \\
        2 & 1 & 2 & 3 & 4 & Invertida \\
        \bottomrule
    \end{tabular}
    \legend{Fonte: elaborado pela autora.}
\end{table}

O escore total é a soma das 8 respostas, de 8 a 32 pontos, e define o
perfil de investidor conforme a Tabela \ref{tab:ipa-classificacao}. Como
a pergunta 1 é sorteada, a análise estatística do Capítulo
\ref{cap:resultados} usa também um escore restrito às perguntas 2 a 8 (de
7 a 28 pontos), que reflete apenas o que o agente efetivamente respondeu.

\begin{table}[htb]
    \caption{Classificação do perfil de investidor por faixa de escore}
    \label{tab:ipa-classificacao}
    \centering
    \begin{tabular}{lcc}
        \toprule
        Perfil & Faixa (soma, 8 perguntas) & Faixa (média por pergunta) \\
        \midrule
        Conservador & 8 a 15  & 1,0 a 1,9 \\
        Moderado    & 16 a 23 & 2,0 a 2,9 \\
        Arrojado    & 24 a 32 & 3,0 a 4,0 \\
        \bottomrule
    \end{tabular}
    \legend{Fonte: elaborado pela autora.}
\end{table}

\subsection{Justificativa da escolha dos instrumentos}

A escolha dos instrumentos deste trabalho, e não de outras versões do
Big Five ou de outros questionários de perfil de investidor, foi
motivada por três razões principais, sendo a terceira delas responsável
pela evolução do instrumento de autorrelato entre o estudo piloto e a
análise principal (seção \ref{sec:bfi44}).

Em primeiro lugar, e mais importante, o IPA e a categorização dos
traços do Big Five em três níveis são os mesmos de
\citeonline{debortoli2019}, o artigo-base cujo instrumento de perfil de
investidor e cuja categorização de traços são adotados diretamente neste
trabalho. Isso vale tanto para o estudo piloto quanto para a análise
principal: a troca do instrumento de \textit{autorrelato} de
personalidade (do BFI-10 para o BFI-44) não alterou essa base comum, já
que o IPA e a categorização Baixo/Médio/Alto (o nível Médio corresponde
ao nível Moderado daquele artigo) permaneceram idênticos nas duas fases. Usar o mesmo IPA que o artigo de referência permite uma
comparação metodologicamente mais direta entre os achados deste trabalho
e os daquele estudo (Capítulo \ref{cap:discussao}).

Em segundo lugar, o IPA é um instrumento de \textit{suitability}
efetivamente utilizado no mercado financeiro brasileiro (adaptado do
questionário de Análise de Perfil do Investidor do Banco do Brasil,
alinhado à ANBIMA e à IOSCO), o que dá ao instrumento,
independentemente dos resultados deste trabalho, relevância prática para
além do contexto acadêmico (seção \ref{sec:justificativa}).

Em terceiro lugar, a escolha do instrumento de \textit{autorrelato} de
personalidade (o que muda entre piloto e análise principal) refletiu um
equilíbrio diferente em cada fase entre custo computacional e granularidade
da checagem de manipulação. O BFI-10, de 10 itens e o mesmo usado por
\citeonline{debortoli2019}, foi o instrumento fixo durante a fase de
comparação entre quatro modelos de linguagem (seção
\ref{sec:selecao-modelo}): com quatro modelos a testar, um instrumento
mais curto reduzia o custo computacional de gerar respostas e
simplificava o \textit{parsing} automático necessário para converter as
respostas textuais em escores numéricos (seção \ref{sec:pipeline}). Uma
vez concluída essa comparação e definido o \texttt{llama3.1:8b} como
modelo da análise principal, esse motivo para preferir um instrumento
curto deixou de se aplicar, já que a análise principal não precisaria
mais ser repetida para vários modelos de qualidade ainda desconhecida. A granularidade adicional do
BFI-44 (mais itens por traço, seção \ref{sec:bfi44}) passou então a
valer o custo computacional extra, o que motivou a troca do instrumento
de autorrelato na análise principal.

\section{\textit{Pipeline} de geração dos agentes}
\label{sec:pipeline}

A geração dos agentes é feita pelo \textit{script} \texttt{src/geracao/gerar\_agentes.py}, disponível no repositório do trabalho (seção \ref{sec:analise-estatistica}). Para cada agente,
o pipeline executa:

\begin{enumerate}
    \item Sorteia um nível (Baixo/Médio/Alto) para cada um dos cinco traços do
    Big Five, de forma independente e com igual probabilidade;
    \item Monta o \textit{system prompt} combinando as descrições
    correspondentes (seção \ref{sec:prompts});
    \item O agente responde ao questionário de autorrelato de personalidade;
    \item O agente responde ao IPA (perguntas 2 a 8; a pergunta 1 é
    atribuída aleatoriamente);
    \item Os escores de ambos os questionários são calculados e a linha
    correspondente é salva em um arquivo CSV.
\end{enumerate}

A Figura \ref{fig:pipeline} resume esse fluxo e as etapas seguintes, da
geração dos agentes à análise estatística e à análise de robustez e
generalização.

\begin{figure}[htbp]
    \caption{\textit{Pipeline} de geração e análise dos agentes}
    \label{fig:pipeline}
    \centering
    \begin{tikzpicture}[
        node distance=6mm,
        etapa/.style={rectangle, rounded corners=2pt, draw=azulfig, fill=azulfig!8,
            line width=0.6pt, align=center, inner sep=5pt, font=\footnotesize,
            text width=#1, execute at begin node={\hyphenpenalty=10000}},
        etapa/.default=12.6cm,
        llm/.style={etapa=5.9cm, fill=azulfig!18, minimum height=2.1cm},
        seta/.style={-{Stealth[length=2.2mm]}, line width=0.7pt, draw=azulescuro}
    ]
        \node[etapa] (sorteio) {\textbf{1. Sorteio do perfil de personalidade}\\
            Um nível (Baixo, Médio ou Alto) para cada um dos cinco traços do Big Five,
            de forma independente e equiprovável; a pergunta 1 do IPA (faixa
            etária) também é sorteada};
        \node[etapa, below=of sorteio] (prompt) {\textbf{2. \textit{System prompt}}\\
            Descrições correspondentes aos níveis sorteados (Quadros
            \ref{quadro:perfis-nivel} e \ref{quadro:prompt-sistema})};
        \node[llm, below=9mm of prompt.south west, anchor=north west] (bfi)
            {\textbf{3a. Chamada ao LLM: BFI-44}\\
            44 itens, escala de 1 a 5\\[2pt]
            $\downarrow$\\
            Escore de autorrelato por traço\\
            (checagem de manipulação)};
        \node[llm, below=9mm of prompt.south east, anchor=north east] (ipa)
            {\textbf{3b. Chamada ao LLM: IPA}\\
            Perguntas 2 a 8\\[2pt]
            $\downarrow$\\
            Escore (Q2--Q8) e perfil de investidor\\
            (Conservador, Moderado ou Arrojado)};
        \node[etapa, anchor=north] (csv) at ($(prompt.south |- bfi.south) + (0,-9mm)$)
            {\textbf{4. Registro}\\
            Respostas convertidas em escores pelo código (até 3 tentativas por
            chamada); uma linha por agente no arquivo CSV, repetida para os
            N=300 agentes};
        \node[etapa, below=of csv] (analise) {\textbf{5. Análise estatística}\\
            Correlação de Spearman e regressão OLS entre o nível induzido e o
            escore do IPA (seção \ref{sec:analise-estatistica})};
        \node[etapa, below=of analise] (robustez) {\textbf{6. Robustez e generalização}\\
            Repetição com outras temperaturas, outro sorteio de perfis e outros
            modelos de linguagem (seção \ref{sec:metodo-robustez})};

        \draw[seta] (sorteio) -- (prompt);
        \draw[seta] (prompt.south) -- ++(0,-4.5mm) -| (bfi.north);
        \draw[seta] (prompt.south) -- ++(0,-4.5mm) -| (ipa.north);
        \draw[seta] (bfi.south) -- (bfi.south |- csv.north);
        \draw[seta] (ipa.south) -- (ipa.south |- csv.north);
        \draw[seta] (csv) -- (analise);
        \draw[seta] (analise) -- (robustez);
    \end{tikzpicture}
    \legend{Fonte: elaborado pela autora. As duas chamadas ao modelo usam o mesmo \textit{system prompt} e não compartilham histórico de conversa.}
\end{figure}

\subsection{Prompts enviados ao modelo}
\label{sec:prompts}

No estudo piloto e nos Experimentos A e B, cada agente corresponde a
duas chamadas independentes ao modelo, feitas pela interface de
programação (API) do Ollama: uma com o questionário de autorrelato e
outra com o IPA. As duas chamadas usam o mesmo \textit{system prompt},
que define a personalidade do agente, e não compartilham histórico de
conversa entre si. Em ambas, a temperatura é 0,7 e a resposta é limitada
a 512 \textit{tokens}; quando a resposta não pode ser interpretada
automaticamente, são feitas até 3 tentativas no total. Os ajustes
adotados nas rodadas da análise de robustez e generalização estão
descritos na seção \ref{sec:metodo-robustez}.

O Quadro \ref{quadro:prompt-sistema} mostra o modelo do \textit{system
prompt}. Os trechos entre chaves são preenchidos pelo código com o nível
sorteado de cada traço e a descrição correspondente do Quadro
\ref{quadro:perfis-nivel}. O Quadro \ref{quadro:prompt-exemplo} mostra o
\textit{system prompt} efetivamente enviado a um dos agentes do
Experimento A. Os nomes dos traços aparecem no \textit{prompt} exatamente
como no código; nele, Amabilidade está escrita como ``Afabilidade'',
sinônimo usado nos arquivos do projeto.

\begin{quadro}[htbp]
    \caption{Modelo do \textit{system prompt} que induz a personalidade}
    \label{quadro:prompt-sistema}
    \begin{promptbox}
    Você deve responder como uma pessoa com o seguinte perfil de personalidade:\\
    - Extroversão (\campo{nível}): \campo{descrição do nível}\\
    - Afabilidade (\campo{nível}): \campo{descrição do nível}\\
    - Consciência (\campo{nível}): \campo{descrição do nível}\\
    - Neuroticismo (\campo{nível}): \campo{descrição do nível}\\
    - Abertura à Experiência (\campo{nível}): \campo{descrição do nível}

    Responda TODAS as perguntas a seguir como essa pessoa responderia, de forma consistente com esse perfil. Não explique, não justifique, apenas dê as respostas no formato pedido.
    \end{promptbox}
    \legend{Fonte: elaborado pela autora. \campo{nível}: Baixo, Médio ou Alto, conforme o sorteio.}
\end{quadro}

\begin{quadro}[htbp]
    \caption{Exemplo de \textit{system prompt} enviado a um agente do Experimento A}
    \label{quadro:prompt-exemplo}
    \begin{promptbox}
    Você deve responder como uma pessoa com o seguinte perfil de personalidade:\\
    - Extroversão (Baixo): É uma pessoa reservada, prefere ambientes tranquilos e evita ser o centro das atenções.\\
    - Afabilidade (Médio): Confia nas pessoas moderadamente: às vezes é crítica, às vezes é tolerante.\\
    - Consciência (Baixo): Tende a ser preguiçosa, desorganizada e costuma deixar tarefas inacabadas.\\
    - Neuroticismo (Médio): Mantém o equilíbrio emocional na maior parte do tempo, mas ocasionalmente fica ansiosa diante de problemas.\\
    - Abertura à Experiência (Médio): Tem alguns interesses artísticos e está moderadamente aberta a novas ideias e experiências.

    Responda TODAS as perguntas a seguir como essa pessoa responderia, de forma consistente com esse perfil. Não explique, não justifique, apenas dê as respostas no formato pedido.
    \end{promptbox}
    \legend{Fonte: dados da pesquisa (primeiro agente do Experimento A).}
\end{quadro}

Os Quadros \ref{quadro:prompt-bfi44} e \ref{quadro:prompt-ipa} mostram as
instruções enviadas em cada chamada, logo após o \textit{system prompt}.
Os dois pedem que o agente escreva o texto exato da alternativa
escolhida, em um formato fixo de uma resposta por linha, que o código
então converte em pontuação. Os itens do BFI-44 e as perguntas
do IPA, abreviados nos quadros, são os dos Quadros
\ref{quadro:bfi44-itens} e \ref{quadro:ipa-perguntas}; a pergunta 1 do
IPA não é enviada ao modelo. No estudo piloto, o \textit{system prompt}
foi o mesmo; mudaram apenas as instruções dos questionários, adaptadas ao
BFI-10.

\begin{quadro}[htbp]
    \caption{Instrução enviada ao modelo para o BFI-44}
    \label{quadro:prompt-bfi44}
    \begin{promptbox}
    Para cada afirmação abaixo, escreva o texto exato da opção escolhida logo após o símbolo $\rightarrow$.\\
    Opções: \textquotedbl Discordo totalmente\textquotedbl{} | \textquotedbl Discordo em parte\textquotedbl{} | \textquotedbl Nem concordo nem discordo\textquotedbl{} | \textquotedbl Concordo em parte\textquotedbl{} | \textquotedbl Concordo totalmente\textquotedbl

    1. Eu me vejo como alguém que Gosta de conversar, é comunicativa. $\rightarrow$\\
    2. Eu me vejo como alguém que Tende a criticar os outros. $\rightarrow$\\
    {[\dots]} \campo{itens 3 a 43, no mesmo formato}\\
    44. Eu me vejo como alguém que É sofisticada em arte, música ou literatura. $\rightarrow$
    \end{promptbox}
    \legend{Fonte: elaborado pela autora.}
\end{quadro}

\begin{quadro}[htbp]
    \caption{Instrução enviada ao modelo para o IPA}
    \label{quadro:prompt-ipa}
    \begin{promptbox}
    Para cada pergunta abaixo, escolha a opção que melhor representa você.\\
    Escreva o texto exato da opção escolhida, sem alterar nada.

    2. Se você realizasse um investimento financeiro hoje (ex: fundos de investimento, poupança, CDI), qual o percentual do total de seu patrimônio que você aplicaria em investimentos?\\
    \hspace*{1.5em}- Até 25\%\\
    \hspace*{1.5em}- Entre 26\% e 50\%\\
    \hspace*{1.5em}- Entre 51\% e 75\%\\
    \hspace*{1.5em}- Mais de 75\%

    {[\dots]} \campo{perguntas 3 a 8, no mesmo formato}

    Responda com exatamente 7 linhas (uma por pergunta, na ordem acima, perguntas 2 a 8). Em cada linha, escreva APENAS o texto exato da opção escolhida, sem letras, sem números, sem qualquer outro texto.\\
    Exemplo (primeiras 2 linhas):\\
    Entre 26\% e 50\%\\
    No máximo 5 anos

    Resposta:
    \end{promptbox}
    \legend{Fonte: elaborado pela autora.}
\end{quadro}

\subsection{Desenho experimental: sorteio aleatório em vez de fatorial}
\label{sec:desenho-fatorial}

Com 5 dimensões $\times$ 3 níveis, um desenho fatorial completo teria
$3^5 = 243$ combinações, cuja repetição um número razoável de vezes seria
inviável no tempo disponível. Optou-se por sortear cada traço de forma
independente por agente, o que permite tratar o nível induzido como
variável ordinal (1/2/3) e usar todos os agentes em uma única análise de
correlação e regressão sobre os cinco traços simultaneamente. Esse é o
desenho do estudo piloto e do Experimento A (seção \ref{sec:exp-a-b}); o
Experimento B usa uma variação desse sorteio em que apenas um traço,
também escolhido por sorteio, varia por agente. A seção \ref{sec:resultados-sorteio} verifica que o sorteio
ficou equilibrado entre os níveis e independente entre os traços (maior
correlação entre pares de traços: $|\rho|$ = \aleatRhoMax), de modo que o
efeito de um traço não se confunde com o de outro.

\subsection{Variáveis do estudo}

\begin{itemize}
    \item \textbf{Variável independente principal}: nível de personalidade
    induzido (Baixo=1, Médio=2, Alto=3) em cada um dos cinco traços,
    conhecido por construção (definido pelo sorteio, sem erro de medida);
    \item \textbf{Variável de checagem de manipulação}: escore do
    autorrelato de personalidade respondido pelo próprio agente;
    \item \textbf{Variável dependente}: escore do IPA restrito às perguntas
    2 a 8 (7--28), usado nas análises estatísticas, e perfil
    (Conservador/\allowbreak Moderado/\allowbreak Arrojado), definido pelo
    escore total (8--32).
\end{itemize}

O nível induzido, e não o autorrelato, foi escolhido como variável
independente principal porque o autorrelato de um LLM pode ser ruidoso ou
inconsistente com a personalidade pedida no \textit{prompt}; usar o nível
induzido evita que esse ruído mascare uma eventual relação real entre
personalidade e perfil de investidor.

\section{Seleção do modelo de linguagem}
\label{sec:selecao-modelo}

Antes de fixar o modelo para a análise principal, quatro modelos de
linguagem foram testados quanto à qualidade da indução de personalidade,
todos com o mesmo instrumento de checagem (BFI-10) e o mesmo
\textit{system prompt}: \texttt{llama3.2:1b} (1,2 bilhão de parâmetros),
\texttt{llama3.2:3b} (3,2 bilhões), \texttt{gemma3:4b} (4,3 bilhões) e
\texttt{llama3.1:8b} (8,0 bilhões). Os três primeiros foram executados
localmente, em unidade central de processamento (CPU); o
\texttt{llama3.1:8b}, por exigir mais capacidade computacional, foi
executado no Google Colab com unidade de processamento gráfico (GPU)
NVIDIA T4.

A Tabela \ref{tab:comparacao-modelos} resume os critérios de comparação.
O critério decisivo foi a checagem de manipulação: o \texttt{llama3.2:1b}
não incorporou nenhum dos cinco traços de forma estatisticamente
detectável e foi descartado; o \texttt{llama3.2:3b} incorporou de forma
significativa quatro dos cinco traços, mas não Abertura à Experiência
($\rho=0{,}077$, n.s.), e com Extroversão fraca ($\rho=0{,}185$); o \texttt{gemma3:4b}, testado apenas em teste reduzido
(N=30), incorporou os cinco traços, mas apresentou viés de
conservadorismo acentuado (60\% dos agentes classificados como
Conservador, 0\% como Arrojado), com uma distribuição de perfil IPA pouco
variável para fins de análise estatística; e o \texttt{llama3.1:8b} foi o
único modelo que reuniu, com N suficiente para inferência, manipulação
significativa nos cinco traços (forte em quatro deles e moderada em
Extroversão), o maior poder explicativo ($R^2$) e a maior variância no
escore do IPA, pré-requisitos para que os efeitos de
personalidade sobre o IPA, se existentes, tenham chance de ser detectados
com significância estatística.

\begin{table}[htb]
    \caption{Comparação entre os quatro modelos de linguagem testados}
    \label{tab:comparacao-modelos}
    \centering
    \footnotesize
    \setlength{\tabcolsep}{4pt}
    \begin{tabular}{>{\raggedright\arraybackslash}p{3.1cm}>{\centering\arraybackslash}p{2.4cm}>{\centering\arraybackslash}p{2.4cm}>{\centering\arraybackslash}p{2.6cm}>{\centering\arraybackslash}p{2.6cm}}
        \toprule
        Critério & \texttt{llama3.2:1b} & \texttt{llama3.2:3b} & \texttt{gemma3:4b} & \texttt{llama3.1:8b} \\
        \midrule
        Manipulação (5 traços)     & Falha total    & Parcial (4/5)  & Boa (5/5, N=30) & \textbf{Significativa (5/5)} \\
        Manipulação (Abertura)     & $-$0,014 n.s.  & 0,077 n.s.    & 0,629***        & \textbf{0,739***} \\
        \% Arrojado no IPA         & 1\%            & 9\%            & 0\%             & \textbf{17\%} \\
        Desvio padrão IPA (Q2--Q8) & 1,96           & 1,99           & 1,64            & \textbf{2,74} \\
        $R^2$ do modelo (OLS)         & --             & 0,067          & -- (N=30)       & \textbf{0,169} \\
        Traços associados ao IPA (Spearman) & Extroversão & Abertura & Extroversão (N=30) & \textbf{Abertura + Extroversão} \\
        Viés de resposta           & Centralidade   & Centralidade   & Conservadorismo & \textbf{Centralidade (atenuado)} \\
        Tempo médio por agente     & 40 s (CPU)     & 49 s (CPU)     & 55 s (CPU)      & \textbf{3,2 s (GPU)} \\
        N analisado                & 316            & 299            & 30              & \textbf{300} \\
        \bottomrule
    \end{tabular}
    \legend{Fonte: dados da pesquisa. Todos os valores de manipulação e $R^2$ referem-se à seleção do modelo, conduzida com o BFI-10 (seção \ref{sec:bfi10-piloto}).}
\end{table}

A progressão do 1b para o 3b e deste para o 8b mostra melhora dentro da
própria família Llama, sobretudo do 3b para o 8b (manipulação, variância
do IPA e poder explicativo), o que sugere que o
tamanho do modelo, e não apenas a família, é um fator relevante para a
validade deste tipo de simulação (discutido em maior profundidade no
Capítulo \ref{cap:discussao}). O \texttt{llama3.1:8b} foi, portanto,
adotado como modelo da análise principal (Experimentos A e B, com o
BFI-44, seção \ref{sec:exp-a-b}); outros modelos aparecem apenas na
análise de robustez e generalização (seção \ref{sec:metodo-robustez}).

\section{Limitações conhecidas do \textit{pipeline}}
\label{sec:limitacoes-pipeline}

\begin{itemize}
    \item \textbf{Pergunta 1 do IPA (faixa etária)}: em testes iniciais,
    100\% dos agentes respondiam a mesma alternativa (viés sistemático);
    por isso, essa pergunta passou a ser atribuída aleatoriamente em vez de
    respondida pelo modelo;
    \item \textbf{Itens de pontuação reversa}: o modelo tende a discordar
    de itens de pontuação reversa (ex.: ``é reservado(a)'') independentemente
    da personalidade induzida, atenuando a manipulação de Extroversão e
    Abertura à Experiência em modelos menores (ver seção
    \ref{sec:selecao-modelo});
    \item \textbf{Temperatura de 0,7 na análise principal}: introduz
    variabilidade entre execuções, ou seja, dois agentes com o mesmo nível
    de personalidade induzido podem gerar respostas diferentes aos
    questionários. O parâmetro \texttt{--seed} do
    \textit{script} de geração controla apenas o sorteio dos níveis de
    personalidade de cada agente (tornando reprodutível \textit{quais}
    combinações de traços foram geradas), não a geração de texto do LLM
    em si. O efeito dessa variabilidade sobre os resultados foi avaliado
    repetindo a análise principal com outras temperaturas (seção
    \ref{sec:metodo-robustez}).
\end{itemize}

\section{Análise de robustez e generalização}
\label{sec:metodo-robustez}

Para avaliar se as associações encontradas no Experimento A dependem de
condições específicas da sua geração, o Experimento A foi repetido
variando um fator de cada vez e mantendo todo o restante do desenho
(indução simultânea dos cinco traços, BFI-44, IPA e mesmos
\textit{prompts}), exceto pelos ajustes técnicos de limite de resposta
descritos a seguir. Foram três frentes:

\begin{enumerate}
    \item \textbf{Sensibilidade à temperatura}: o Experimento A, gerado com
    temperatura 0,7, foi repetido com temperaturas 0,0, 0,3 e 1,0, com o
    mesmo modelo (\texttt{llama3.1:8b}) e N=300 em cada rodada. Todas as
    rodadas usaram a mesma semente do sorteio (\texttt{--seed} 123), de
    modo que os 300 perfis de personalidade e a pergunta 1 do IPA sorteada
    são idênticos aos do Experimento A, agente a agente; a única diferença
    entre as rodadas é a temperatura;
    \item \textbf{Novo sorteio de perfis}: o Experimento A foi repetido com
    outra semente (\texttt{--seed} 456), mesmo modelo, temperatura 0,7 e
    N=300, o que produz 300 perfis de personalidade sorteados de novo. Esse
    teste avalia se os resultados dependem do sorteio específico usado na
    análise principal. O novo sorteio passou por verificações de
    equilíbrio e independência análogas às da seção
    \ref{sec:resultados-sorteio};
    \item \textbf{Outros modelos de linguagem}: o Experimento A foi
    repetido com outros modelos, mantendo temperatura 0,7 e a semente 123
    (mesmos perfis do Experimento A, agente a agente). Três modelos
    abertos, executados via Ollama no Google Colab com GPU T4, foram
    rodados com N=300: \texttt{qwen2.5:7b}, \texttt{mistral:7b} e
    \texttt{gemma3:12b}. Para esses modelos, o limite de resposta foi
    ampliado de 512 para 1.536 \textit{tokens}, porque alguns deles
    repetem o texto de cada item do BFI-44 antes da resposta. Dois modelos
    maiores da família Nemotron (NVIDIA), o Nemotron 3.5 Super e o
    Nemotron 3 Ultra, foram acessados pela API do OpenRouter em sua cota
    gratuita, que limita o número de chamadas por dia; por isso, foram
    rodados apenas como teste exploratório, com N=30 (os mesmos perfis dos
    30 primeiros agentes do Experimento A), com o raciocínio interno do
    modelo desligado e limite de resposta de 4.096 \textit{tokens}. Um terceiro
    modelo previsto para essa etapa, o Gemma 4 31B, não pôde ser incluído:
    o provedor que o disponibiliza na cota gratuita recusou quase todas as
    chamadas por congestionamento ao longo de vários dias de tentativas,
    e apenas 2 dos 30 agentes foram gerados, número insuficiente para
    qualquer análise.
\end{enumerate}

Além das análises da seção \ref{sec:analise-estatistica}, as rodadas de
temperatura e de outros modelos, por usarem os mesmos perfis do
Experimento A, permitem uma comparação pareada, agente a agente: a correlação de Spearman entre o escore IPA obtido na
rodada e o obtido no Experimento A, e a porcentagem de agentes que
receberam o mesmo perfil de investidor nas duas rodadas.

Complementarmente, a qualidade da indução no Experimento A foi avaliada
também pela \textbf{correlação entre os traços autorrelatados}: como os
níveis induzidos são sorteados de forma independente (seção
\ref{sec:resultados-sorteio}), uma correlação relevante entre os escores
de autorrelato de dois traços pode ser atribuída ao próprio modelo, e não
ao sorteio. Foram
calculadas a correlação de Spearman entre os escores autorrelatados dos
10 pares de traços e a correlação entre o nível induzido de cada traço e
o autorrelato de cada um dos outros quatro, que indica o quanto a
indução de um traço ``vaza'' para outro.

\section{Análise estatística}
\label{sec:analise-estatistica}

Dois métodos estatísticos foram aplicados a cada um dos três conjuntos de
dados principais deste trabalho (estudo piloto, Experimento A e
Experimento B), sempre sobre os agentes com perfil IPA válido
(Conservador, Moderado ou Arrojado):

\begin{enumerate}
    \item \textbf{Correlação de Spearman} ($\rho$) entre o nível de
    personalidade induzido (1/2/3) e o escore do IPA (soma das perguntas
    2 a 8, variando de 7 a 28), calculada separadamente para cada um dos
    cinco traços. A correlação de Spearman, não paramétrica, foi preferida
    à de Pearson por não pressupor relação linear nem normalidade das
    variáveis, o que a torna adequada a uma variável independente ordinal
    de apenas três níveis;
    \item \textbf{Regressão linear múltipla por mínimos quadrados
    ordinários (OLS)}, com o escore do IPA como variável dependente e os
    cinco níveis de personalidade induzidos (como variáveis numéricas
    1/2/3) como regressores simultâneos, replicando a abordagem adotada
    por \citeonline{jiangpengyan2024}. Ao contrário da correlação de
    Spearman por traço, a regressão OLS estima o efeito de cada traço
    \textit{controlando} o efeito dos demais, o que permite distinguir um
    efeito próprio de um traço de um efeito decorrente apenas de sua
    correlação com outro traço mais relevante.
\end{enumerate}

Os dois métodos foram aplicados adicionalmente ao par (nível induzido,
escore do autorrelato de personalidade) de cada traço, para a checagem de
manipulação (seção \ref{sec:bfi44}; resultados no Capítulo
\ref{cap:resultados}). No Experimento B, a mesma correlação de Spearman
foi calculada também de forma estratificada, apenas entre os agentes em
que aquele traço específico era o ``traço dominante'' (seção
\ref{sec:exp-a-b}), já que a maioria dos agentes tem cada traço
individual fixo em ``Médio'' e, portanto, sem variância a correlacionar.

O nível de significância adotado em todo o trabalho é $\alpha=0{,}05$;
resultados são reportados com três níveis de significância convencionais
(* $p<0{,}05$; ** $p<0{,}01$; *** $p<0{,}001$; ``n.s.'' para $p \geq 0{,}05$). As
análises foram implementadas em Python, com as bibliotecas
\texttt{pandas} (manipulação dos dados), \texttt{scipy.stats}
e \texttt{statsmodels}. As mesmas
análises foram aplicadas, sem nenhuma alteração, a cada rodada da análise
de robustez e generalização (seção \ref{sec:metodo-robustez}), sempre
sobre o escore IPA das perguntas 2 a 8.

O código, os instrumentos e os dados gerados neste trabalho estão
disponíveis em
\url{https://github.com/Marianaferrarez/tcc-bigfive-perfil-investidor}.

\section{Uso de ferramentas de inteligência artificial}
\label{sec:uso-ia}

Ferramentas de inteligência artificial generativa foram utilizadas como
apoio em três etapas deste trabalho: na automatização das consultas às
APIs dos modelos de linguagem, de modo a aproveitar as janelas de tempo e
os limites de uso das cotas gratuitas; no refinamento dos \textit{prompts}
usados para induzir a personalidade e aplicar os questionários; e no apoio
à escrita da monografia, principalmente na formatação e na resolução de
problemas do documento em \LaTeX. Todo o
material produzido com esse apoio foi revisado pela autora, que se
responsabiliza integralmente pelo conteúdo do trabalho.
